%% kapa: exam=1 round=27 year=2016 period=1 form=A subject=civil_law subject_ko=민법 seq=14 number=14
%% source: Q-Net 감정평가사 기출문제 (cst00309) artlSeq=5169402 file=extracted/1008/5169402/제27회 감정평가사 1차 문제/제27회 감정평가사 1교시문제(A형).pdf
%% conversion: kordoc 4.24.51 -> tools/gen1.py (auto); review=proofread-sonnet; flags=none
\begin{question}{14}
법률행위의 조건에 관한 설명으로 옳은 것은?
\begin{choices}
  \item 조건의 성취가 미정한 권리라도 일반규정에 의하여 처분하거나 상속할 수 있다.
  \item 조건의 성취로 인하여 불이익을 받을 당사자가 신의성실에 반하여 조건의 성취를 방해하였어도 상대방은 그 조건이 성취한 것으로 주장할 수 없다.
  \item 모든 법률행위에는 조건을 붙일 수 있다.
  \item 조건이 법률행위 당시 이미 성취한 것인 경우에는 그 조건이 정지조건이면 그 법률행위는 무효로 한다.
  \item 조건이 법률행위 당시 이미 성취할 수 없는 것인 경우에는 그 조건이 정지조건이면 조건 없는 법률행위로 한다.
\end{choices}
\end{question}
\answer{1}
